\section{Geometry, motives, and the Grothendieck programme}
\label{sec:geometry}

The letters A and M have the dependency meaning fixed in the introduction.
In particular, accepting the stipulated endpoint A is not a certification of
every intermediate statement in the note, and assuming M does not certify
the functional analysis in its proposed proof. The findings below are new
formulations proved in this run; no claim of priority over the literature is
intended. Earlier findings are used with their original identifiers.

\subsection{The transfer results already established}
\label{ag:transfer}

\paragraph{AG-AT (the exact Artin--Tate domain).}
For a finite-type scheme over $\mathbb Z$ with exact all-prime point counts
\[
 \#X(\mathbb F_{p^r})=
 \sum_j p^{jr}\operatorname{Tr}(F_p^r\mid V_j^{I_p}),
\]
where the $V_j$ are finitely many virtual finite-image representations,
the identity
\[
 \zeta_X(s)=\prod_j L(V_j,s-j)
\]
is unconditional. This is ZETA\_TWO ZT01, in
\texttt{RETURN\_ZETA\_TWO\_final\_IMPLICATIONS.md}, Section 1,
extending BRIDGES F1. If the representations have abelian image, A gives
noninteger global divisor support on the lines $\Re s=j+1/2$.
For arbitrary finite image, M gives the corresponding meromorphic
divisor-support assertion through integer Brauer induction (ZT02).
Negative exponents can produce poles; Artin holomorphy is an extra question.

\paragraph{AG-DED and AG-LOCAL (ramification and models).}
The permutation representation on the embeddings of a number field has
inertia invariants giving the Euler factors of $\operatorname{Spec}\mathcal
O_K$ at ramified primes as well. Thus every Dedekind zeta function is covered
by M, and abelian fields over $\mathbb Q$ are covered by A (ZT03).
If a factorization holds only away from a finite set $S$, the actual factors
at $S$ must be restored (ZT04). The counterexample ZT05,
\[
 \zeta_{\operatorname{Spec}\mathbb Z[1/p]}(s)
       =(1-p^{-s})\zeta(s),
\]
has additional zeros $2\pi i k/\log p$, $k\ne0$. Such local divisor lines
are not shifted global critical lines. This is a statement about the
integral model, not a failure of the generic-fibre motivic normalization.

\paragraph{AG-CM (the potentially CM domain).}
ZETA\_TWO CM01 supplies the all-weight extension to pure objects in the
absolute-Hodge category generated by Artin motives, Tate motives, and
potentially CM abelian varieties. Its cited CM descent input gives a signed
product of Hecke functions of one weight; norm twisting makes these
characters unitary. The required M therefore includes infinite-order
unitary characters. The result is meromorphic divisor RH, and does not
declare absolute-Hodge projectors algebraic. The exact prior locator is
\texttt{RETURN\_ZETA\_TWO\_final\_IMPLICATIONS.md}, Section 2, CM01,
with its reported Anderson Section 5.7, Corollary 5.7.2 and Remark 5.7.3.
The original source has now been recovered and checked: G.~W.~Anderson,
\emph{Cyclotomy and an extension of the Taniyama group}, Compositio
Math.~57 (1986), 153--217. The category is defined in Section 5.3,
printed p.~178 (PDF p.~27); Theorem 6, Section 5.6, printed pp.~179--180
(PDF pp.~28--29), identifies its one-dimensional characters with algebraic
Hecke characters, including the inertia and Frobenius comparisons.
Corollary 5.7.2 and Remark 5.7.3 are both on printed p.~181 (PDF p.~30).
The former gives integer exponents, with no positivity claim; the latter
gives the common weight. Its internal reference to ``Corollary 5.8.2'' is
a printed cross-reference error, not a new result used here.

\paragraph{AG-CM-REC and AG-CM-PERIOD (classical conclusions already
stronger than zero location).}
In this precise potentially CM absolute-Hodge category, equality of the
whole $L$-functions implies isomorphism of objects: Anderson,
Proposition 5.7.5, printed pp.~181--182 (PDF pp.~30--31). This is
unconditional. The proof uses semisimplicity and Corollary 5.7.4's
identification, for a weight-zero object, of the pole order at $1$ with
the dimension of its invariant part. Equality of the local factors makes
the tensor products with a dual agree away from finitely many primes;
their local factors are nonzero and nonpolar at $1$. A pole of
$L(M\otimes M^\vee,s)$ therefore gives a nonzero morphism
$M\longrightarrow N$, and induction on the common rank proves the
claim. The proof was read in full, rather than inferred from GRH.
This is a recognition theorem for semisimple motives; it supplies no
effective finite bound for general infinite CM monodromy.

Anderson's Section 5.8, Theorem 7, printed p.~182 (PDF p.~31), also
proves Deligne's period conclusion
$L(M,0)/c^+(M)\in\mathbb Q$ for the stated \emph{tractible, critical}
objects: those for which the $L$-function is a single algebraic Hecke
$L$-function over a totally real or CM field. This is a separate
classical special-value theorem, using Siegel and Blasius, and is not
enlarged to arbitrary motives by A or M. It complements, rather than
replaces, the cycle/height and integral comparison gates below.

\paragraph{AG-WEIGHT and AG-NONCM (the residual geometric factors).}
With geometric Frobenius,
\[
 \operatorname{wt}N(n)=\operatorname{wt}N-2n,
 \qquad L(N(n),s)=L(N,s+n).
\]
The line for weight $w$ is $\Re s=(w+1)/2$. Positive weight is not the
remaining obstruction: CM objects of positive weight are covered by M,
whereas $\operatorname{Sym}^2h^1(E)(1)$ for a non-CM elliptic curve still
has a non-CM factor despite having weight zero. ZT06 reduces a finite-type
scheme to finitely many smooth projective geometric factors and actual
exceptional fibres. ZT07 proves that, under A, an off-line elliptic zero
in $1/2<\Re s<3/2$ cannot be cancelled by the two Tate factors in the
good-model formula
\[
 \zeta_{\mathcal E}(s)
   =\zeta^S(s)\zeta^S(s-1)/L^S(E,s).
\]
These are retained reductions, not a proof of RH for the non-CM factor.
Framework locators carried from ZETA\_TWO are SGA 4 XVII Sections 5.1--5.3;
SGA 5 XV Section 3, no.~2, unnumbered theorem, equation (1);
EGA IV$_3$ (8.8.2); and Deligne, Weil I, Theorem 1.6, and Weil II,
Theorem 3.3.1.

\subsection{AG-FROB: a finite recognition theorem for Artin--Tate objects}
\label{ag:frob}

Let $L/\mathbb Q$ be a nontrivial finite Galois extension, with group $G$,
absolute discriminant $D_L$, and degree $n_L$. Put
\[
 B_L=(4\log D_L+\tfrac52 n_L+5)^2.
\]
Let
\[
 U=\bigoplus_j V_j(-j),\qquad W=\bigoplus_j W_j(-j),
\]
where all sums are finite and $V_j,W_j$ are actual finite-dimensional
characteristic-zero representations of $G$. Their Euler polynomials at
unramified primes are understood as those of the associated compatible
Artin--Tate systems; a single auxiliary prime may instead be chosen greater
than $B_L$.

\textbf{Statement.} Assume RH for $\zeta_L$ and $\zeta_{\mathbb Q}$.
If
\[
 \det(1-TF_p\mid U)=\det(1-TF_p\mid W)
 \quad\text{for every }p\le B_L\text{ unramified in }L,
\]
then $V_j\simeq W_j$ for every $j$, and hence $U\simeq W$.
For rational coefficients this is also an isomorphism in the rational
Artin--Tate category. A supplies the hypothesis when $L/\mathbb Q$ is
abelian; M supplies it in general. When $L=\mathbb Q$, the one prime $p=2$
suffices.

\begin{proof}
At an unramified prime, every eigenvalue in $V_j(-j)$ has the form
$p^j\omega$, where $\omega$ is a root of unity. Therefore its complex
absolute value is $p^j$. Distinct integers $j$ give distinct absolute
values, so equality of the total Euler polynomials separates uniquely into
equality of the polynomials in each weight. It follows that
$\operatorname{Tr}(F_p\mid V_j)=\operatorname{Tr}(F_p\mid W_j)$ for every
$j$ and every prime in the hypothesis.

Bach--Sorenson, \emph{Explicit bounds for primes in residue classes},
Math. Comp. 65 (1996), Theorem 5.1, printed p.~1729 (PDF p.~13), gives
an unramified prime $p\le B_L$ in every conjugacy class of $G$.
The precise analytic assumptions are stated in their Section 2,
printed p.~1720: zero-freeness of $\zeta_L$ and $\zeta_{\mathbb Q}$ in
$\Re s>1/2$. The reduction from cyclic to general Galois extensions is
Corollary 3.3, p.~1721. Thus the characters of $V_j$ and $W_j$ agree on
every conjugacy class. In characteristic zero finite-group representations
are semisimple, and equal characters give isomorphic representations.

For completeness, this last assertion also descends to a non-algebraically
closed characteristic-zero coefficient field $E$. The space of
$G$-equivariant maps is the solution space over $E$ of finitely many linear
equations. If it contains an invertible map after scalar extension, its
determinant polynomial is nonzero on that $E$-vector space. Since $E$ is
infinite it has a nonzero value at an $E$-point. That point is an invertible
$G$-equivariant map over $E$. For $L=\mathbb Q$, the eigenvalues $2^j$
already recover all multiplicities from the polynomial at $2$.
\end{proof}

In Grothendieck's language this is effective recognition inside a specified
semisimple Artin--Tate realization category. It extends
ALPOGE\_BRIDGES AQ-FINITE, whose sharper abelian conductor bound remains
available. It does not construct algebraic projectors for an arbitrary
geometric realization. Moreover, recognition of an Artin motive is weaker
than recognition of a finite cover as a scheme; the connected Gassmann
example in CX-GEO-004's companion finding makes this distinction explicit.
The forbidden-prime variant is handled there by a quadratic compositum.

\paragraph{AG-COVER: equal Artin motives need not give equal covers.}
Put $L=\mathbb Q(\sqrt2,\sqrt3)$ and let $K_1,K_2,K_3$ be its three
quadratic subfields. The finite schemes
\[
 X=\operatorname{Spec}L\amalg\operatorname{Spec}\mathbb Q
       \amalg\operatorname{Spec}\mathbb Q,
 \qquad
 Y=\coprod_{i=1}^3\operatorname{Spec}K_i
\]
are not isomorphic, but their rational Artin motives, and all their local
Artin Euler factors, are equal. Consequently
\[
 \zeta_L(s)\zeta_{\mathbb Q}(s)^2
       =\zeta_{K_1}(s)\zeta_{K_2}(s)\zeta_{K_3}(s).
\]
This identity includes the factors at ramified primes.

\begin{proof}
Write $G=\operatorname{Gal}(L/\mathbb Q)=C_2\times C_2$, with its three
subgroups $H_i$ of order two. The character of
$\mathbb Q[G]\oplus\mathbb Q^2$ takes value $6$ at the identity and
$2$ at every other element. Each nonidentity element belongs to exactly
one $H_i$, so the character of $\bigoplus_i\mathbb Q[G/H_i]$ has the
same values. Characteristic-zero semisimplicity and the descent argument
in AG-FROB give an isomorphism of rational representations. Isomorphic
representations have isomorphic inertia invariants with their residual
Frobenius action, proving equality of every local factor and hence the
zeta identity in the initial half-plane and by continuation. On the other
hand the degrees of the connected components are $4,1,1$ for $X$ and
$2,2,2$ for $Y$, precluding a scheme isomorphism.
\end{proof}

The same example over $\mathbb Z$ uses the respective rings of integers;
over $\mathbb Z[1/6]$ these are finite \emph{\'etale} covers. Thus even
complete Euler data forget the multiplication on an Artin algebra. A
applies to all fields here because they are abelian over $\mathbb Q$,
but does not restore the forgotten multiplication. The independent
transitive seven-element example CX-GEO-005 treats the connected
$G$-set distinction without asserting an arithmetic realization.

\subsection{AG-LAMBDA: the precise closure supplied by tensor operations}
\label{ag:lambda}

\textbf{Statement.} The finite direct sums of abelian Artin characters with
integer Tate twists are closed under tensor products, duals, exterior and
symmetric powers, and direct summands in characteristic zero. All their
$L$-functions have the A divisor-support conclusion. Replacing abelian
finite image by arbitrary finite image gives the corresponding M conclusion.
The same operations preserve the potentially CM absolute-Hodge category
used in CM01, so its M conclusion applies there as well.

\begin{proof}
Over an algebraic closure of the coefficient field an abelian Artin--Tate
object is a sum of lines $\chi_a(-j_a)$. A tensor product or a monomial
appearing in a symmetric or exterior power is a line
$(\prod_a\chi_a^{e_a})(-\sum_a e_aj_a)$. The character still has finite
abelian image and the Tate index is an integer. Duality replaces the
character and index by their inverses and negatives. Direct summands remain
in the same semisimple category, and descent preserves the asserted
factorization. The nonabelian assertion follows because all these
operations still factor through a finite group and commute with Tate
twists; apply ZT02 to their resulting Artin representations. Finally, the
potentially CM category is the tensor, dual, and summand closure in CM01.
Symmetric and exterior powers are the images of the usual idempotents in
the symmetric-group action on tensor powers, so belong to it too.
\end{proof}

The placement is SGA 6 V Section 2 and Section 7, especially (7.1.1),
and the determinant/$\lambda$-ring discussion carried in ZETA\_TWO ZT08
and SV01. These closures do not generate an arbitrary non-CM $GL_2$
motive from CM objects; enlarging a list of tensor operations is not the
missing analytic theorem.

\subsection{AG-SIGN: central signs with the necessary normalization}
\label{ag:sign}

\textbf{Statement.} Let $F$ be a nonzero holomorphic function on a
neighbourhood of $[c,\infty)$, real-valued there. Suppose $F$ has no zeros
on $(c,\infty)$ and is positive at one sufficiently large real point.
If $r=\operatorname{ord}_{s=c}F$, then
\[
 \frac{F^{(r)}(c)}{r!}>0.
\]
If, moreover, $F(c+z)=\varepsilon F(c-z)$ near $0$, with
$\varepsilon\in\{1,-1\}$, then $(-1)^r=\varepsilon$.
For a real, standardly normalized, holomorphic Hecke or CM $L$-function,
the absence of real zeros to the right of the centre follows from M
together with its usual analytic continuation and Euler-product theory.
The analogous Dirichlet statement uses A whenever no pole is present.

\begin{proof}
A continuous real nonvanishing function has constant sign on the connected
interval $(c,\infty)$. The normalization makes that sign positive.
Write $F(c+t)=a_rt^r+O(t^{r+1})$, with $a_r\ne0$. For $t>0$ the
left-hand side is positive, so $a_r>0$ after division by $t^r$ and passage
to the limit. Comparing the leading terms in the functional equation gives
$a_r=\varepsilon(-1)^r a_r$, hence the parity assertion.
\end{proof}

\paragraph{AG-HEIGHT (what this proves for a supplied height formula).}
If $F(c)=0$, then $F'(c)\ge0$: it is positive when $r=1$ and zero when
$r>1$. Consequently, an independently established identity
$H(Z)=C F'(c)$ with a specified constant $C>0$ yields $H(Z)\ge0$.
The elementary deduction is complete; the existence, domain, and
normalization of the height identity are extra inputs. It does not assert
that every homologically trivial cycle has such a formula, or that its
$L$-function has simple central vanishing. For an identity involving the
first nonzero derivative, the sign is strictly positive. For a formula with
an additional sign, that sign must be retained.

\paragraph{AG-SIGN-NORM (a counterexample to the unnormalized inference).}
Zero locations and functional-equation parity alone do not determine the
sign: $F(s)=s-c$ and $-F(s)$ have exactly the same zeros and the same odd
functional equation, but their central derivatives have opposite signs.
This is a counterexample to the stated analytic inference, not an example
of two standardly normalized arithmetic $L$-functions.

These findings refine \texttt{OURS\_HODGE\_BRIDGE\_SKETCH\_20261121.md},
Section 3.5. Their archive placement is cote 44, scans 2, 4 and 6,
Theorems 4--5, and scans 8--9, the letter to Serre of 12 August 1964,
as reported in that attached sketch. No original scan has been inferred
from this sign argument. The remaining arithmetic Hodge-index gate is
the cycle/height comparison for the stated class of cycles, not a sign
convention suppressed in a formula.

\subsection{AG-CRITICAL: why off-centre critical integers gain little}
\label{ag:critical}

\textbf{Statement.} Suppose $N$ is a pure, integral-weight, explicitly
Hecke-factored motive of weight $w$, with its usual functional equation
and archimedean factors. At an endpoint of the unitary strip, require each
factor in the specified signed Hecke factorization (or its dual, when the
functional equation is used) to be nonpolar. This factorwise requirement
prevents a denominator pole from producing a zero. Let $m$ be critical in
Deligne's sense: $L_\infty(N,s)$ and $L_\infty(N^\vee,1-s)$
are both regular at $s=m$. If $m\ne(w+1)/2$, then the standard
zero-free half-plane and boundary for unitary Hecke $L$-functions,
together with the functional equation, already give $L(N,m)\ne0$.
M is not needed for this conclusion.

The precise unitary boundary theorem used here is
D.~Ramakrishnan--R.~J.~Valenza, \emph{Fourier Analysis on Number Fields}
(1999), Section 7.7, Theorem 7-28, printed p.~289: for every number
field and every unitary idele class character, there is no zero at
$1+it$. Section 7.4, Theorem 7-19, printed p.~279, supplies the completed
continuation and functional equation, with its pure norm-character pole
exceptions retained. These are classical inputs independent of M.

\begin{proof}
After unitary normalization set $z=m-w/2$. If $w$ is even, $z$ is an
integer; if $w$ is odd, $z$ is a half-integer. An integer $m$ with
$0<z<1$ therefore has $z=1/2$, and is precisely the excluded centre.
For every other $m$ one has $z\ge1$ or $z\le0$. In the first case
standard unitary Hecke nonvanishing in $\Re z\ge1$, away from a pole,
gives the result for each factor and hence its finite signed product.
In the second case the functional equation relates the completed value
to its dual at $1-z\ge1$. Criticality ensures that both archimedean
factors are finite; gamma functions have no zeros. The nonzero dual
value therefore gives a nonzero original value. The factorwise nonpolar
condition excludes a Tate pole or its inverse being misread as a finite
nonzero special value. For a product with nonnegative exponents, the
usual no-pole condition implies this condition.
\end{proof}

Thus the strongest new real-axis information from M lies between the
off-centre integer points and at signs of central leading coefficients;
it does not evaluate regulators, periods, central ranks, or integral
torsion. This is the special-value distinction in ZETA\_TWO Section 4
and in the Hodge sketch's edge-versus-centre table. The relevant motivic
conventions are Deligne, \emph{Valeurs de fonctions $L$ et p\'eriodes
d'int\'egrales}, Sections 1.1--1.3; the period conjectures numbered 1.8
and 2.8 are separate inputs.

\subsection{AG-DET: Hecke versions of the square-variable determinant}
\label{ag:det}

The following formulation records exactly which analytic data are needed.
Let $\Lambda$ be a nonzero entire function of order at most one, with
\[
 \Lambda(c+z)=\varepsilon\Lambda(c-z),\qquad
 \varepsilon\in\{1,-1\},
\]
and let $r=\operatorname{ord}_{s=c}\Lambda$. Assume that every other zero
is $c+i\gamma$, with $\gamma$ real and nonzero. Choose one positive
$\gamma$ from each pair, with multiplicity, and write the sequence as
$\gamma_1,\gamma_2,\ldots$ in nondecreasing order, so $\gamma_1$
denotes the least positive zero when the sequence is nonempty.
Define the removable entire function
\[
 H(t)=t^{-r}\Lambda(c+it),\qquad
 D(w)=H(\sqrt w)/H(0).
\]
No choice of square-root branch affects $D$.

\textbf{Statement.} The diagonal operator
$B=\operatorname{diag}(\gamma_j^{-2})$ is positive and trace class, and
\[
 D(w)=\prod_j(1-w\gamma_j^{-2})=\det(I-wB).
\]
In particular this applies under M to standard entire selfdual completed
Hecke functions and to the entire completed $L$-function of a CM elliptic
curve, with its centre $c=1$. The usual order, functional equation,
and completion are additional classical inputs, not assertions supplied
by the word ``GRH''. Poles must first be removed in a separately specified
entire completion; arbitrary meromorphic Artin quotients are not included.

\begin{proof}
Comparison of the first nonzero Taylor coefficients shows
$\varepsilon=(-1)^r$. Hence $H(-t)=H(t)$, so $D$ is entire by the even
power series for $H$. Its order is at most one half: a bound of order
$1+\delta$ for $H(t)$ gives a bound of order $(1+\delta)/2$ for $D(w)$.
Its zeros, with multiplicity, are exactly $\gamma_j^2$.
The convergence exponent of the zeros of an entire function of order at
most one is at most one. In particular
$\sum_j|\gamma_j|^{-2}<\infty$. Equivalently the diagonal positive
operator $B$ has finite trace norm. Hadamard factorization for $D$, whose
order is strictly less than one, has no nonconstant exponential factor
and uses genus-zero factors. Since $D(0)=1$, its constant factor is one.
This gives the displayed product. The Fredholm determinant of a diagonal
trace-class operator is that same absolutely locally uniformly convergent
product, proving the identity. An exact analytic locator for the
factorization used here is C.~McMullen, \emph{Advanced Complex Analysis},
Math 213a course notes, version 16 November 2026, Theorem 3.16,
printed p.~92 (PDF p.~93); Theorems 3.10--3.12 give the zero-exponent
bound used just above. This is classical background, not a new claim
dependent on that revision date.
\end{proof}

This extends the prior zeta realization BR-TR-01 to a precisely specified
Hecke domain. Its relationship to cote 1 scans 79--80 and 111--112 is
the trace/Fredholm theorem reported in
\texttt{OURS\_INTAKE\_RESEARCH\_SA23\_20261120.md}, Section 4.1,
geo-H-lidskii and geo-H-genus. Those sheets are not themselves an
arithmetic RH claim. The operator here is defined from the zeros; the
construction has not supplied a cohomology theory or an operator defined
independently from prime data.

\subsection{AG-MOMENT and AG-HANKEL: recoverable positive spectral data}
\label{ag:moment}

Assume the hypotheses of AG-DET and infinitely many distinct positive
zeros. Put $\lambda_j=\gamma_j^{-2}$ and
\[
 s_k=\operatorname{Tr}(B^k)=\sum_j\lambda_j^k\quad(k\ge1).
\]
The logarithmic determinant expansion is
\[
 -\log D(w)=\sum_{k\ge1}\frac{s_k}{k}w^k
 \qquad (|w|<\gamma_1^2).
\]

\textbf{AG-MOMENT.} For every $k\ge1$,
\[
 s_k^{-1/k}\le\gamma_1^2\le\frac{s_k}{s_{k+1}}.
\]
The left bounds increase to $\gamma_1^2$; the right bounds decrease to it.
If $m_1$ is the multiplicity of the least positive zero, then
$\lim_{k\to\infty}s_k\gamma_1^{2k}=m_1$.

\begin{proof}
Let $a=\max_j\lambda_j=\gamma_1^{-2}$. Then
$s_k\ge a^k$ and $s_{k+1}\le a s_k$, proving the two inequalities.
The $\ell^p$ norm monotonicity gives
$s_{k+1}^{1/(k+1)}\le s_k^{1/k}$. One direct proof is to divide by
$s_k^{1/k}$: the resulting nonnegative numbers have $k$-th powers
summing to one, are each at most one, and thus have $(k+1)$-th powers
summing to at most one. Hence the left bounds increase.
Cauchy--Schwarz applied to
$\lambda_j^{k/2}$ and $\lambda_j^{(k+2)/2}$ gives
$s_{k+1}^2\le s_ks_{k+2}$, so the right bounds decrease.

There are finitely many indices with $\lambda_j=a$, since $\sum\lambda_j$
is finite. For $k\ge1$,
$0\le(\lambda_j/a)^k\le\lambda_j/a$, whose sum is finite.
Dominated convergence shows
$s_k/a^k\to m_1$. It follows that both $s_k^{1/k}$ and
$s_{k+1}/s_k$ tend to $a$, which proves both limits and the
multiplicity formula.
\end{proof}

\textbf{AG-HANKEL.} For all integers $h\ge1$ and $d\ge0$, the matrix
\[
 (s_{h+i+j})_{0\le i,j\le d}
\]
is positive definite. Its determinant is therefore strictly positive.

\begin{proof}
For a nonzero vector $a=(a_0,\ldots,a_d)$ and
$P(X)=\sum_{i=0}^d a_iX^i$, its Hermitian quadratic form equals
\[
 \sum_{j\ge1}\lambda_j^h|P(\lambda_j)|^2.
\]
This converges because the $\lambda_j$ are bounded and
$\sum\lambda_j^h<\infty$. All terms are nonnegative. A nonzero
polynomial has only finitely many distinct zeros, whereas the positive
$\lambda_j$ have infinitely many distinct values. At least one term
is positive. This proves strict positive definiteness.
\end{proof}

These consequences strengthen the two-by-two moment inequalities carried
in BR-TR-02. They describe actual determinant coefficients available in
the cote 1 analytic framework, while leaving the arithmetic/geometric
realization gate unchanged. They assert no simplicity and no
random-matrix statistics.

\subsection{AG-SIGNED: a finite moment witness for a surviving pole}
\label{ag:signed}

Consider a finite signed product of specified entire selfdual functions
of order at most one, with a common centre $c$, each satisfying the
zero-location hypotheses of AG-DET. Let $r\in\mathbb Z$ be the net
central order. Divide out the central Laurent factor and normalize its
nonzero constant term. Its square-variable quotient has the form
\[
 D(w)=\prod_{\lambda\in\mathcal L}(1-w\lambda)^{m_\lambda},
 \qquad m_\lambda\in\mathbb Z\setminus\{0\},\quad
 \lambda>0,\quad \sum_\lambda |m_\lambda|\lambda<\infty,
 \qquad D(0)=1.                                      \tag{AG-S1}
\]
Here $\mathcal L$ has no positive accumulation point, and the signed
multiplicities are taken \emph{after} cancellation. The form (AG-S1)
follows by applying AG-DET to each factor. More generally the proposition
below applies whenever (AG-S1) is independently supplied. Put
\[
 s_k=\sum_\lambda m_\lambda\lambda^k\quad(k\ge1),\qquad
 H_d=(s_{1+i+j})_{0\le i,j\le d}.
\]
These moments are the local logarithmic coefficients
$-\log D(w)=\sum_{k\ge1}s_kw^k/k$ near $0$.

\textbf{Statement.} All the matrices $H_d$ are positive semidefinite
if and only if $m_\lambda>0$ for every $\lambda\in\mathcal L$.
Equivalently, the normalized quotient has no noncentral poles. The
central order condition $r\ge0$ must be checked separately. If a
negative multiplicity survives, a finite matrix $H_d$ already detects it,
with the explicit witness described below.

\begin{proof}
Set $a=\sup\mathcal L$ and introduce the finite signed measure
\[
 \mu=\sum_\lambda m_\lambda\lambda\,\delta_\lambda
 \quad\hbox{on }[0,a].
\]
The empty support case is immediate. For every complex polynomial
$P(X)=\sum_{i=0}^d b_iX^i$, absolute convergence gives
\[
 \sum_{i,j=0}^d\overline b_i s_{1+i+j}b_j
       =\int |P(x)|^2\,d\mu(x).
\]
Positive multiplicities therefore imply all the required inequalities.

Conversely, suppose $m_{\lambda_0}<0$. The point $\lambda_0>0$ is
isolated in the compact set $\{0\}\cup\mathcal L$. Let
\[
 \delta=\operatorname{dist}
   (\lambda_0,(\{0\}\cup\mathcal L)\setminus\{\lambda_0\})>0,
 \qquad R>a+\lambda_0,
 \qquad q=1-\delta^2/R^2<1.
\]
For each nonnegative integer $N$, the polynomial
\[
 P_N(x)=\left(1-\frac{(x-\lambda_0)^2}{R^2}\right)^N
\]
has degree at most $2N$, takes the value $1$ at $\lambda_0$, and has
$|P_N(\lambda)|\le q^N$ at every other point of $\mathcal L$.
Writing $W=\sum_{\lambda\ne\lambda_0}|m_\lambda|\lambda$, we obtain
\[
 \int |P_N|^2\,d\mu
   \le -|m_{\lambda_0}|\lambda_0+Wq^{2N}.
\]
Choose $N$ with $Wq^{2N}<|m_{\lambda_0}|\lambda_0$; take $N=0$
if $W=0$. Then the coefficient vector of $P_N$ is a negative
quadratic-form witness for the matrix $H_{2N}$, of size $2N+1$.
This proves the converse and its finite-witness assertion. In (AG-S1)
a negative net multiplicity is exactly a pole at $w=\lambda^{-1}$,
which proves the analytic equivalence. Removing the central Laurent
factor erased its order, explaining the separate condition on $r$.
\end{proof}

Geometrically this strengthens the cote 1 determinant bridge in the
meromorphic setting. A or M can first localize an identified divisor;
positivity of these finite moment matrices is then an exact additional
condition preventing cancellation from leaving poles. It is not supplied
by zero location alone. In particular, universal Hecke GRH plus Brauer
induction does not prove this positivity, and hence does not prove Artin
holomorphy. The criterion is for the specified quotient: symmetrizing
an arbitrary nonselfdual function with its dual can cancel a pole between
the two factors, so holomorphy of that product must not be attributed to
each unsymmetrized factor. The local character-matrix test in MT-ART02
is a complementary criterion at an individual Artin divisor point.

\subsection{AG-EXT: the precise information that the determinant forgets}
\label{ag:ext}

The earlier findings BR-MIX-01 and ZETA\_TWO SV01--SV02 are retained.
At a good prime an adapted matrix for a mixed object is block triangular,
so its determinant equals that of its weight-graded object. The nonsplit
Kummer extension from the unit $2$ over $\mathbb Z[1/2]$ and its split
Tate object therefore have the same incomplete $L$-function. The omission
of ramified factors is essential: inertia invariants need not preserve a
short exact sequence. These findings refute recovery of extensions from
semisimplified good Euler data; they do not assert that all richer
arithmetic data agree.

The exact independent background locator is Deligne--Goncharov,
\emph{Groupes fondamentaux motiviques de Tate mixte},
arXiv:math/0302267v2 (27 February 2003), (2.1.3), printed p.~14
(PDF p.~16):
\[
 \operatorname{Ext}^1(\mathbb Q(0),\mathbb Q(n))=
 \begin{cases}
 0&n\le0,\\
 \mathcal O_S^\times\otimes\mathbb Q&n=1,\\
 K_{2n-1}(k)\otimes\mathbb Q&n\ge2.
 \end{cases}
\]
Propositions 2.2--2.3 on the same page identify the corresponding graded
free motivic Lie algebra, while Section 2.4 explicitly distinguishes a
free presentation from a canonical set of generators. These statements
are unconditional background. They do not calculate integral torsion or
the pro-$\ell$ kernels in the prompt's cote 149 programme.

\subsection{AG-KUMMER: an integral obstruction on the actual cote 149 sheets}
\label{ag:kummer}

The original Montpellier PDF was inspected here, not merely the intake.
It has a cover page: PDF p.~80 has repository footer 79, and so on.
At PDF p.~80/footer 79, equations (13)--(17) relate an inertia-normalizer
extension, its $H^2$ class, and Kummer cohomology. PDF p.~81/footer 80,
(18)--(19), identifies the relevant line bundle as the normal bundle of
the section. PDF p.~82/footer 81, (20)--(21), and PDF p.~85/footer 84,
(22), describe the $H^1$ and completed-unit ambiguity. These are the
actual geometric objects on the cited sheets. Their link to LM6
A.44--A.45 is supplied by ORDER item 31; no theorem about an irregular
prime or $K_8(\mathbb Z)$ is printed on these pages.

Here is a precise integral consequence, independent of A and M. Let $S$
be a scheme on which $\ell$ is invertible, put $P=\operatorname{Pic}(S)$,
and let $\operatorname{Br}'(S)=H^2_{\text{ét}}(S,\mathbb G_m)$.
Write $\widehat P_\ell=\varprojlim_nP/\ell^nP$ and
$T_\ell A=\varprojlim_n A[\ell^n]$, with transition maps multiplication
by $\ell$. There are exact sequences
\begin{align*}
 0&\longrightarrow\widehat P_\ell
   \longrightarrow\varprojlim_nH^2_{\text{ét}}(S,\mu_{\ell^n})
   \longrightarrow T_\ell\operatorname{Br}'(S)\longrightarrow0,\\
 0&\longrightarrow\Gamma(S,\mathcal O_S^*)^{\wedge}_\ell
   \longrightarrow\varprojlim_nH^1_{\text{ét}}(S,\mu_{\ell^n})
   \longrightarrow T_\ell P\longrightarrow0.
\end{align*}
We deliberately use inverse-limit cohomology: identifying it with a
chosen definition of continuous cohomology requires its own
$\varprojlim{}^1$ check. If $P$ is finitely generated, the second right
term is zero. A line bundle $\mathcal L$ has zero compatible integral
$\ell$-adic Chern class if and only if its class in $P$ is torsion of
order prime to $\ell$.

\begin{proof}
The Kummer exact sequence gives
\[
 0\longrightarrow P/\ell^nP
 \longrightarrow H^2(S,\mu_{\ell^n})
 \longrightarrow\operatorname{Br}'(S)[\ell^n]\longrightarrow0,
\]
and the analogous degree-one sequence with units modulo $\ell^n$th
powers and $P[\ell^n]$. The exact source is Stacks Project,
Lemma 59.28.1, tag 03PK, and its immediately following cohomology
sequence. The transition maps on the left terms are surjective.
Consequently a compatible sequence on the right lifts compatibly in
the middle: after choosing a lift at level $n$, correct any lift at
level $n+1$ by a lift of its discrepancy in the left term. This recursive
construction proves right-surjectivity after taking inverse limits;
left-exactness is immediate. It proves both displayed sequences.

Now write the finitely generated group $P$ as
$\mathbb Z^a\oplus P[\ell^\infty]\oplus P_{\ell'}$, with finite torsion
summands. Its $\ell$-adic completion is
$\mathbb Z_\ell^a\oplus P[\ell^\infty]$, and the kernel of completion
is exactly $P_{\ell'}$. Also no nonzero compatible infinitely
$\ell$-divisible torsion sequence can occur in a finitely generated
group, so $T_\ell P=0$. The injection in degree two then proves the
Chern-class criterion.
\end{proof}

\paragraph{AG-KUMMER-2: a fully explicit integral/rational separation.}
Let $K=\mathbb Q(\sqrt{-35})$, $\omega=(1+\sqrt{-35})/2$, and
$S=\operatorname{Spec}\mathbb Z[\omega,1/2]$. The invertible ideal
$\mathfrak p=(3,\omega)$ determines a line bundle $\mathcal L$ whose
compatible integral $2$-adic Chern class is nonzero, although its image
after tensoring with $\mathbb Q_2$ is zero.

\begin{proof}
The ring of integers is $\mathbb Z[\omega]$ and its discriminant is
$-35$. The minimal polynomial $X^2-X+9$ is irreducible modulo $2$,
so $(2)$ is the only prime over $2$, and is principal. Modulo $3$ it is
$X(X-1)$; thus the primes over $3$ are
$\mathfrak p=(3,\omega)$ and $\overline{\mathfrak p}=(3,\omega-1)$.
Since $N(\omega)=9$ and $\omega$ is a unit modulo
$\overline{\mathfrak p}$, unique ideal factorization gives
$\mathfrak p^2=(\omega)$. The ideal $\mathfrak p$ is not principal:
an integral generator would have norm $3$, whereas
\[
 N(a+b\omega)=a^2+ab+9b^2=(a+b/2)^2+35b^2/4
\]
cannot be $3$ for integers $a,b$.

Minkowski's bound is $2\sqrt{35}/\pi<4$, so every ideal class has a
representative of norm $1$, $2$, or $3$. There is no ideal of norm $2$;
the two of norm $3$ have the same class of order two. Hence the class
group is $\mathbb Z/2$. The precise classical bound is Milne,
\emph{Algebraic Number Theory}, v3.08 (19 July 2020), Theorem 4.3,
PDF p.~72. Inverting $2$ quotients the class group by the classes of
the primes over $2$, which here are principal, so
$\operatorname{Pic}(S)=\mathbb Z/2$ as well. To see this last assertion
directly, fractional ideals of the localization are obtained by forgetting
the valuation at $(2)$; a relation introduced in the localization differs
from a principal ideal only by a power of that principal prime.

AG-KUMMER now injects this nonzero class into the inverse limit of
$H^2(S,\mu_{2^n})$. It has order two there, so rationalization kills it.
The total space of $\mathcal L$ has a zero section with normal bundle
$\mathcal L$, making the example a literal normal-bundle obstruction of
the kind appearing in cote 149. No claim identifying this particular
scheme's group-cohomology extension class is needed.
\end{proof}

This proves a concrete limit of rational information: a rationally zero
normal Chern class need not be integrally zero. A supplies Dirichlet GRH
for the quadratic field here, but changes none of this calculation.
The finite-level Kummer test is therefore a genuine extra gate for an
integral inertia-splitting argument. It does not refute the original
LM6 statements, whose comparison maps must still be supplied.

\subsection{AG-C34-TYPE and AG-C34-REPAIR: resolving the commutator notation}
\label{ag:c34repair}

The official cote 34 PDF has a cover: PDF p.~53 is manuscript/footer
p.~52, and PDF p.~64 is p.~63. Theorem 1.1 on the former page fixes a
field $k$ finitely generated over its prime field, a normal geometrically
connected $k$-scheme $X$ locally of finite type, a prime
$\ell\ne\operatorname{char}k$, and an open subgroup
$\pi\subseteq\pi_1(X)$. It defines $\pi_0$ to be the image in
$G_k$ and $I=\pi\cap\pi_1(X_{\bar k})$, in
\[
 1\longrightarrow I\longrightarrow\pi\longrightarrow\pi_0
 \longrightarrow1.
\]
Its $I(\ell)$ is the maximal abelian pro-$\ell$ quotient of $I$.
The theorem states that $I(\ell)$ is finitely generated over
$\mathbb Z_\ell$ and that its $\pi_0$-coinvariants are finite.
The proof was located and read in Sections 2--3, PDF pp.~70--78
(manuscript pp.~69--77), rather than guessed from the intervening
applications. Section 2.4, PDF p.~72, makes the dual $H^1$-invariant
reduction; Section 2.5, PDF pp.~73--74, uses a curve compactification,
its cyclotomic boundary, and the finite-field curve theorem after
spreading out. Sections 3.1--3.3, PDF pp.~75--78, reduce normal $X$ to
a generic hyperplane curve using surjectivity on fundamental groups.
These are classical geometric inputs; A and M do not enter them.

In the crossed-out Remark 1.17 on PDF p.~64, $H$ is a quotient of $I$
whose kernel is $\pi$-stable, and the readable final hypothesis is that
$H$ is a solvable pro-$\ell$ group. The preceding proposed
$\ell$-adic Lie-group language is itself struck out. The remark writes
\[
 1\longrightarrow H\longrightarrow\pi'\longrightarrow\pi_0
 \longrightarrow1
\]
and places the closure of $[\pi_0,H]$ as an open subgroup of $\pi'$.
This is a tentative, struck manuscript assertion, not a theorem.

\paragraph{AG-C34-TYPE (the literal ambient-group obstruction).}
For an infinite $\pi_0$, every commutator of a lift of an element of
$\pi_0$ with an element of $H$ lies in $H$, and its closure cannot be
open in $\pi'$. Indeed the quotient by $H$ is the infinite group
$\pi_0$, whereas a subgroup open in a profinite group has finite index.
For a concrete example satisfying the manuscript's geometric hypotheses,
take $X=\mathbb G_m/\mathbb F_q$ and $\ell\nmid q$. Its geometric
pro-$\ell$ quotient $H=\mathbb Z_\ell(1)$ is abelian and hence
solvable; the rational point $1$ gives a section. Arithmetic Frobenius
acts as multiplication by $q$, so the action commutators generate
$(q-1)\mathbb Z_\ell(1)$. This is open in $H$ but is not open in
$\pi'$, whose quotient is $G_{\mathbb F_q}\simeq\widehat{\mathbb Z}$.
Thus the literal ambient group in the inherited Q03 cannot be retained.
For nonabelian $H$ there is a second issue: $\pi_0$ acts canonically
only by outer automorphisms unless a section or lifts are specified.

\paragraph{AG-C34-REPAIR (the intrinsic corrected statement).}
Let $H$ be \emph{any} pro-$\ell$ quotient of $I$ with $\pi$-stable
kernel, and let $P$ be the resulting quotient extension of $\pi_0$ by
$H$. Define $K$ to be the closed subgroup generated by all
$[p,h]$, $p\in P$, $h\in H$. Then $K$ is open in $H$, and
\[
 H/K\simeq (H^{\mathrm{ab}})_{\pi_0},
 \qquad [H:K]\le |I(\ell)_{\pi_0}|<\infty.       \tag{AG-C34.1}
\]
In particular the solvability and Lie-group hypotheses are unnecessary
for this precise corrected conclusion.

\begin{proof}
Because $H$ is normal in $P$, conjugating a commutator $[p,h]$ by
$a\in P$ gives $[apa^{-1},aha^{-1}]$. Hence $K$ is normal in $P$.
It contains the closed commutator subgroup of $H$, so $H/K$ is abelian;
all conjugation by $P$ is trivial on it. Conjugation on
$H^{\mathrm{ab}}$ factors canonically through $\pi_0$, irrespective
of any section. Killing its action differences is exactly the additional
quotient by the images of the commutators $[p,h]$. This proves the
isomorphism in (AG-C34.1).

The surjection $I\twoheadrightarrow H$ induces an equivariant
surjection $I(\ell)\twoheadrightarrow H^{\mathrm{ab}}$ by the
universal property of the maximal abelian pro-$\ell$ quotient.
Taking coinvariants remains surjective: it is simply quotienting both
compact abelian groups by the closed subgroups generated by their
action differences. Theorem 1.1 therefore supplies a surjection from
the finite group $I(\ell)_{\pi_0}$ to $H/K$. This proves the bound.
Since $K$ is closed, finite index in the compact group $H$ implies
openness (the finite Hausdorff quotient is discrete).
\end{proof}

This repairs the archived Q03 as a theorem already implied by its
stated abelian predecessor. It does not identify the unqualified
symbol $[\pi_0,H]$ in every possible intended meaning. A smaller
commutator subgroup formed from a chosen section, or openness in a
different ambient group, needs a separately stated question. In
particular the repair is not described as an automatic proof of all
interpretations of the crossed-out remark.

\paragraph{AG-C34-INDEX (an explicit Frobenius refinement).}
More generally, let $M=H^{\mathrm{ab}}$ be a finitely generated
$\mathbb Z_\ell$-module with continuous $\Gamma$-action, let $T$ be
its finite torsion subgroup, and let $\Lambda=M/T$. Suppose some
$F\in\Gamma$ has no eigenvalue $1$ on
$\Lambda\otimes\mathbb Q_\ell$. Then
\[
 |\operatorname{coker}(F-1\mid M)|
   =|T^F|\ell^{v_\ell\det(1-F\mid\Lambda\otimes\mathbb Q_\ell)},
 \qquad
 |M_\Gamma|\le
 |T^F|\ell^{v_\ell\det(1-F\mid\Lambda\otimes\mathbb Q_\ell)}.
                                                               \tag{AG-C34.2}
\]
Equality in the second assertion holds if $\Gamma$ is topologically
generated by $F$.

\begin{proof}
The assumption gives $\Lambda^F=0$. Apply the snake lemma to
$0\to T\to M\to\Lambda\to0$ with vertical maps $F-1$.
It yields
\[
 0\longrightarrow T/(F-1)T\longrightarrow M/(F-1)M
 \longrightarrow\Lambda/(F-1)\Lambda\longrightarrow0.
\]
For an endomorphism of the finite group $T$, kernel and cokernel have
the same cardinality. Smith normal form over $\mathbb Z_\ell$ gives
the order of the last group as the displayed power of $\ell$.
This proves the equality. The continuous image $(F-1)M$ is compact
and hence closed; the full coinvariant quotient is a quotient of it.
If the powers of $F$ are dense in $\Gamma$, every action difference
is a limit of differences of powers, all of which belong to
$(F-1)M$, so the two quotients agree.
\end{proof}

For a smooth proper geometrically connected curve $C/\mathbb F_q$,
let $H$ be its full geometric pro-$\ell$ fundamental group and $J$ its
Jacobian. Then $H^{\mathrm{ab}}=T_\ell J$, compatibly with Galois
action, and the corrected commutator quotient is
\[
 H/\overline{[P,H]}
       \simeq T_\ell J/(F-1)T_\ell J
       \simeq J(\mathbb F_q)[\ell^\infty].             \tag{AG-C34.3}
\]
For the fundamental-group comparison, an exact modern primary locator
is J.~S.~Milne, \emph{Jacobian Varieties}, Proposition 9.1,
pp.~30--32, together with \emph{Abelian Varieties}, Remarks 15.3
and 15.5, pp.~28--29. To check the last isomorphism directly, use
$0\to T_\ell J\to V_\ell J\to J(\overline{\mathbb F}_q)
[\ell^\infty]\to0$. The fixed subgroup of $T_\ell J$ is zero:
$J(\mathbb F_q)$ is finite, so a compatible indefinitely divisible
sequence of its $\ell$-power torsion points vanishes. Thus $F-1$ is
invertible on $V_\ell J$, and the snake lemma gives (AG-C34.3).
Consequently the index is exactly $\ell^{v_\ell\#J(\mathbb F_q)}$.
This finite-field calculation even uses less than the full Weil bound.
In higher-dimensional smooth proper situations purity supplies the
corresponding no-eigenvalue-$1$ test, but finite-field purity is already
classical and no new global GRH input is involved.

\subsection{Archive and run consequences, with the original scope retained}
\label{ag:archive}

\paragraph{AG-C1.} Cote 1/79--80 and 111--112 gives the trace/Fredholm
framework, and SA23 geo-H-lidskii/geo-H-genus is its reported modern
reconstruction. A yields the earlier zeta positivity realization;
M yields AG-DET for the stated Hecke functions. AG-MOMENT and AG-HANKEL
give further consequences in that framework. A geometric operator defined
from primes remains a separate construction.

\paragraph{AG-C16.} Cote 16, scans 14--18, no.~7, concerns the integral
Lefschetz correspondence ring. Scans 44--50 contain the Coates letter's
conditions A, C and index positivity; the attached Hodge sketch uses a scan
numbering one higher than the GCT\_NSRH intake. These are algebraicity and
integrality problems. Neither A nor M constructs the required
correspondence or bounds its denominator. The sketch's explicit route is
the rank comparison on $X\times X$, not a claim that zero location
reverses Grothendieck's implication to the Weil conjectures.

\paragraph{AG-C34.} Original PDF p.~74, manuscript p.~73, (2.5.5),
has now been visually checked: its multiplicative $1$ denotes the trivial
invariant group, equivalently additive $0$. The finite-field curve theorem
is the stated input, so inherited Q04 is unchanged. The original
definitions and full proof route for Theorem 1.1 have also been read.
AG-C34-TYPE identifies the literal obstruction in Q03, and
AG-C34-REPAIR proves the coherent intrinsic correction. AG-C34-INDEX
gives its finite Frobenius/Jacobian index. No assertion of new purity or
of a consequence of A/M is made.

\paragraph{AG-C44.} Cote 44 Theorems 4--5 and the Serre letter, at the
locators in AG-SIGN, concern polarization and heights. AG-SIGN corrects
the central-derivative bridge but leaves cycle existence and the
height formula explicit. The supplied Hodge sketch also cites a separate
2026 CM-Hodge theorem for its abelian positivity consequences; those
claims are not consequences of either A or M and are carried with
their original dependency rather than independently certified here.

\paragraph{AG-C149 and AG-GT.} ORDER\_tonight, item 31, points to cote
149 scans 79--85, LM6 A.44--A.45, integral bridge B3, and the
irregular-prime/$\ell=2$ questions. GRH does not remove the integral
cohomology comparison or the specified class-group component gate.
The semi-graded Euler information in AG-EXT explains why the endpoint
zero theorems do not evaluate that extension data. In SA22 run 1,
A6-node-character is the quadratic Artin character times the cyclotomic
character; A gives effective detection of that finite quadratic sign.
M additionally covers finite nonabelian splitting fields in effective
Frobenius tests such as AG-FROB. These improvements do not reverse the
recorded refutation of run 2 A:inertia-maximality: finite-index inertia
subgroups have the same weights. Nor do they identify the profinite GT
kernel from the cyclotomic quotient. The existing supplied claim about
rational Deligne--Drinfeld freeness (CC-GT03, external 2026 input) is
carried separately from the full profinite Ihara question R-GTHAT.

\paragraph{AG-PURITY and AG-WM.} Finite-field purity and the usual
Grothendieck--Lefschetz trace formula are already classical. The full
geometric weight--monodromy statement remains represented in these
materials by BR-WM-02's finite exterior-power trace criterion; its
remaining inequalities are geometric hypotheses, not consequences of
global Dirichlet or Hecke zero locations. Likewise the supplied archive
questions about standard conjectures require algebraic cycles or
correspondences. No formal logical independence from A or M is asserted;
the precise point is that the displayed deductions do not supply those
inputs.
